\documentclass[11pt]{article}

\usepackage{amsmath,amssymb}
\usepackage{xurl}
\usepackage[hidelinks]{hyperref}
\setlength{\emergencystretch}{2em}

% Shared commands for notes in docs/paper-gaps/.
%
% Two halves.  The link commands below are project identity: OWNER/REPO is the
% placeholder that scripts/bootstrap_project.py replaces with this project's
% GitHub slug and Pages site.  Everything after them is NOTATION, and notation
% belongs to the source paper: the definitions here are an example set, and a
% project replaces them with the macros of its own paper's style file.  The one
% rule that never changes: a command defined here must mean exactly what the
% source paper's macro means, or a note silently says something else.

\newcommand{\Lean}{\textsc{Lean}}
\newcommand{\leanid}[1]{\textsf{#1}}
\newcommand{\ghissue}[1]{\href{https://github.com/OWNER/REPO/issues/#1}{GitHub issue \##1}}
\newcommand{\ghpr}[1]{\href{https://github.com/OWNER/REPO/pull/#1}{PR \##1}}
% Local issue tracker (issues/NNNN-*.md in this repository); no remote URL.
\newcommand{\localissue}[1]{issue \##1 (\path{issues/})}
\newcommand{\doclink}[1]{\href{https://OWNER.github.io/REPO/docs/#1.html}{\path{#1}}}

% --- Notation: replace with the source paper's own macros --------------------
% \F, \N, \C, \R, \tr, \ot, \eps, \E, \bx, \bu, \bv, \by

\newcommand{\F}{\mathbb{F}}
\newcommand{\N}{\mathbb{N}}
\newcommand{\C}{\mathbb{C}}
\newcommand{\R}{\mathbb{R}}
\newcommand{\tr}{\mathrm{tr}}
\newcommand{\eps}{\epsilon}
\newcommand{\ot}{\otimes}
\newcommand{\E}{\mathop{\bf E\/}}

% bold random variables
\newcommand{\bu}{\boldsymbol{u}}
\newcommand{\bv}{\boldsymbol{v}}
\newcommand{\bx}{{\boldsymbol{x}}}
\newcommand{\by}{\boldsymbol{y}}

% T1 so literal < and > in \texttt placeholders typeset as themselves
\usepackage[T1]{fontenc}

% bra-ket
\usepackage{braket}

% --- Example structured spaces ---
\newcommand{\polyfunc}[3]{\mathcal{P}(#1, #2, #3)}
\newcommand{\polymeas}[3]{\mathrm{PolyMeas}(#1, #2, #3)}
\newcommand{\polysub}[3]{\mathrm{PolySub}(#1, #2, #3)}

% Approximation relations (paper uses \approx_\eps and \simeq_\zeta)
\newcommand{\approxeps}{\approx_{\varepsilon}}
\newcommand{\simgeq}{\simeq_{\zeta}}


% Local notation for this note.
\DeclareMathOperator{\poly}{poly}
\newcommand{\Id}{\mathrm{Id}}
\renewcommand{\ghissue}[1]{%
  \href{https://github.com/Dengnifer/MIPStarRE-A/issues/#1}{GitHub issue \##1}}

\title{The Error Term of the Combined-Lines Consistency Lemma}
\author{}
\date{2026-08-30}

\begin{document}
\maketitle

\section{At a glance}

\begin{itemize}
  \item \textbf{Difficulty: easy.}  No step of the retained proof route is
  in doubt.  The obstruction is that the error function claimed in the
  statement of the combined-lines lemma of the Pauli basis analysis,
  $\poly(m^2\eps, md/q)$, is not delivered by either of the two
  derivations printed in its source proof: one route establishes
  $m \cdot \poly(\eps, md/q)$, the other would give
  $m^2 \cdot \poly(\eps, md/q)$ and rests in addition on a distributional
  decomposition that the stated sub-line lemma does not supply.  The
  correction is to state the first route's established error form and confirm
  that the sole consumer tolerates it.
  \item \textbf{Key Mathlib inputs:} the
  Cauchy--Schwarz inequality for state-dependent operator distances
  (project-local, blueprint Chapter 12), mixture decompositions of
  question distributions, and elementary averaging over restricted
  line distributions; no new Mathlib surface beyond finite
  expectations and inner-product estimates already used by the
  low individual degree development.
  \item \textbf{Estimated weight:} a \textbf{statement-level error-form
  decision} plus about \textbf{one lemma interface} of checking that the
  consumer absorbs an $m$ prefactor into its constants; the three chained
  claims behind the retained route are ordinary Cauchy--Schwarz estimates,
  on the order of \textbf{a few hundred lines} of formal proof in total
  when that part of the development is reached.
  \item \textbf{Mathlib/project split:} essentially \textbf{100\%
  project-local}.  The estimates use only the project's state-dependent
  distance calculus, the sub-line distribution, and elementary
  probability on finite sets.
  \item \textbf{Key inputs:} the combined point and line measurements and
  the sub-line distribution of the Pauli basis analysis in
  \cite{Ji2020MIPStar}, and the error function of the classical soundness
  theorem imported from \cite{Ji2020LowIndividualDegree}, into which the
  corrected error form must be absorbed.  The blueprint correction is tracked
  in \ghissue{173}.
\end{itemize}

\section{Key theorem forms}

Throughout, $(q,m,d)$ is an admissible parameter tuple of the Pauli basis
test, $\ket{\hat\psi}$ is the expanded state of the projective strategy
under analysis, succeeding with probability at least $1-\eps$, and all
approximations are with respect to $\ket{\hat\psi}$.  For lines
$\ell_X, \ell_Z$ in $\F_q^m$, the POVM $\{T^{\ell_X,\ell_Z}_{f_X,f_Z}\}$,
with outcomes pairs of polynomials of degree at most $md$ on the two
lines, is the two-basis combined line measurement, and
$\{\hat Q^{x,z}_{a,b}\}$ is the combined point measurement; their
consistency errors are $\delta_P = \poly(\eps, md/q)$ and
$\delta_Q = \poly(\eps)$ respectively, and
$\delta_{\mathrm{Line}} = \poly(\eps)$ is the line--point consistency
error of the strategy.  The lemma under discussion asserts that a
measurement $\{\hat Q^{\ell}_f\}$, built for each line $\ell$ in
$\F_q^{2m+2}$, is $\delta_{\mathrm{combine}}$-consistent, on average over
the line-point distribution over $\F_q^{2m+2}$, with the combined point
measurement on the extended coordinates, together with the
register-swapped counterpart.%
\footnote{In the mirror of the source consulted here,
\path{references/qpbt-paper/14_analysis_of_the_pauli_basis_test.tex}: the
lemma statement, with the error form $\poly(m^2\eps, md/q)$, is at
lines~1020--1034 (label \path{lem:qld-4-13}); the central display with
the first-route bound at lines~1134--1137
(label \path{eq:qld-combined-lines-consistency}); the three claims at
lines~1140--1239 (labels \path{claim:17-1}--\path{claim:17-3}); the
second-route conclusion $O(m^2\delta_P)$ at lines~1241--1245, preceded by
a duplicated sentence whose second copy carries a malformed reference;
the product-distribution estimate invoked there at lines~1055--1061
(label \path{eq:qld-xz-lines-restricted}); and the sub-line lemma at
lines~1063--1116 (label \path{lem:qld-sublines}).}
The comparison involves four forms of the error function
$\delta_{\mathrm{combine}}(\eps,m,d,q)$.
\begin{enumerate}
  \item \textbf{Source form.}
  $\delta_{\mathrm{combine}} = \poly(m^2\eps,\, md/q)$.  This is the form
  printed in the statement.  Neither printed derivation supports it.
  \item \textbf{Established form (first route).}
  $\delta_{\mathrm{combine}}
  = O\big(m\,(\eps^{1/4} + \delta_P^{1/4} + \delta_Q^{1/4}
  + \delta_{\mathrm{Line}}^{1/2})\big)$, a function of the form
  $m \cdot \poly(\eps,\, md/q)$.
  \item \textbf{Second-route form (conditional).}
  $\delta_{\mathrm{combine}} = O(m^2 \delta_P)
  = m^2 \cdot \poly(\eps,\, md/q)$, conditional on a joint-law
  decomposition that the sub-line lemma, as stated, does not provide.
  \item \textbf{Consumer requirement.}  The unique consumer, the
  polynomial-pair lemma, feeds $\delta_{\mathrm{combine}}$ into an error
  of the form
  $a(md)^a(\delta_{\mathrm{combine}}^b + q^{-b} + 2^{-bmd})$ imported
  from \cite{Ji2020LowIndividualDegree}; any bound
  $m^{O(1)} \cdot \poly(\eps, md/q)$ is absorbed by enlarging $a$ and
  shrinking $b$, so forms~2 and~3 both suffice downstream.
\end{enumerate}

\section{The domain of the conditional Lean helper}

The source statement does not take an arbitrary joint point measurement as an
input.  It uses the family $\{\hat Q^{x,z}_{a,b}\}$ constructed by the preceding
point lemma, whose error has the form $\delta_Q(\eps)=\poly(\eps)$.  Consequently,
an internal theorem which accepts an already constructed point witness must retain
the same quantitative dependency: it should first fix a function
$\delta_Q:\mathbb R\to\mathbb R$ satisfying the one-variable polynomial-error
predicate, then choose $\delta_P$, and at strategy error $\eps$ accept only a
point witness with bound $\delta_Q(\eps)$.

Before the repair tracked by \ghissue{389}, the Lean helper
\path{exists_combinedLinesWitness_ofPointsWitness} instead had the quantifier
order
\[
  \exists\delta_P=\poly(\eps,md/q)\quad
  \forall\eps\ \forall\delta_Q\ \forall Q,
  \qquad Q\text{ has point error at most }\delta_Q
  \Longrightarrow T\text{ has line error at most }\delta_P(\eps,md/q).
\]
Thus the line-error function was chosen before an unrestricted scalar
$\delta_Q$, and did not depend on it.  The fields of
\path{CombinedPointsWitness} bound self-consistency and the two ordered-product
comparisons by precisely this scalar; they impose no relation between an
arbitrary $\delta_Q$ and $\eps$.  The conditioned-pasting estimate therefore
cannot eliminate $\delta_Q$ from its conclusion.

There is a concrete mathematical obstruction to the old domain.  Set
$\eps=0$, take $m=1$ and fixed $d$, and let $q$ grow through admissible field
sizes with $q>d+1$.  On any perfect strategy, define the supplied projective
joint point measurement deterministically, with X answer $x^{d+1}$ and a fixed
Z answer, identically on the two player sides.  Its self-consistency defect is
zero, while the two ordered-product defects are finite and hence can be covered
by a sufficiently large scalar $\delta_Q$, which the old helper quantified
arbitrarily.  On an axis X line, however, every outcome allowed by
\path{CombinedLinesWitness.axis_degree_X} evaluates as a polynomial of degree at
most $d$.  Its difference from $x^{d+1}$ has at most $d+1$ roots, so even after
mixing line outcomes the disagreement probability on that line is at least
$1-(d+1)/q$.  Since axis lines have fixed positive weight in the line-point
law, the resulting line consistency defect stays bounded away from zero as
$q\to\infty$.  In contrast, the two-variable polynomial-error predicate forces
$\delta_P(0,d/q)\to0$.  This argument diagnoses the theorem interface; no Lean
counterexample theorem is asserted here.

The repaired conditional helper has the form
\[
  \forall\delta_Q,\quad \delta_Q=\poly(\eps)\Longrightarrow
  \exists\delta_P=\poly(\eps,md/q)\quad
  \forall\eps\ \forall Q\text{ with error }\delta_Q(\eps),\quad
  \exists T\text{ with error }\delta_P(\eps,md/q).
\]
This is not a new hypothesis on the paper theorem.  The source-facing Lean
declaration \path{exists_combinedLinesWitness} preserves its original statement:
it existentially obtains $\delta_Q$ and its point witnesses from the preceding
point construction, then applies the conditional helper.

The preserved proof route isolates four named obligations.  First,
\path{exists_combinedLine_restored_defect_le} constructs the conditioned XZX
line measurement and retains the numerical point error in its raw defect bound.
Second, \path{exists_conditioned_polynomial_bound} uses the hypothesis
$\delta_Q=\poly(\eps)$ to choose the two-variable line-error function.  Third,
\path{nondegenerateLinePastingMass_bounds} controls the conditioning mass between
$1/2$ and $1$.  Fourth, \path{consistencyDefect_opposite_symm} restores every
directed opposite placement.  After these are available in the current modular
tree, the conditional theorem can reuse their proof from protected commit
\path{0f4ef05370350f4017439ebd839ef0561f13130f}; the source-facing theorem then
combines it with \path{exists_combinedPointsWitness}.  Until those construction
obligations are recovered and checked, both theorem proofs remain explicit proof
obligations rather than bridge assumptions.

At the time of this repair, \path{exists_conditioned_polynomial_bound} is already
available in \path{MIPStarRE/QPBT/Combining/ErrorBounds.lean}.  The remaining
proof extension is a focused modular recovery of the XZX measurement definition,
its two axis-degree lemmas, \path{exists_combinedLine_restored_defect_le},
\path{nondegenerateLinePastingMass_bounds}, and
\path{consistencyDefect_opposite_symm}.  It must reuse the current pasting and
line-measurement APIs rather than importing the protected monolithic file.

\section{The source assertion and the two routes}

The source proof in \cite{Ji2020MIPStar} constructs $\hat Q^{\ell}_f$ by
averaging, over a distribution $D$ on tuples $(\ell, \ell_X, \ell_Z)$
provided by its sub-line lemma, the pushforward of
$T^{\ell_X,\ell_Z}$ along a combining map, and reduces the lemma to a
bound on the scalar consistency defect
\begin{equation}
  \Lambda \,=\,
  \E_{(\ell,\ell_X,\ell_Z) \sim D}\
  \E_{(x,z,\alpha,\beta) \in \ell}
  \sum_{f_X,f_Z} \bra{\hat\psi}
  \big(T^{\ell_X,\ell_Z}_{f_X,f_Z}\big)_{\mathsf{A}\mathsf{A}'} \ot
  \big(\Id - \hat{Q}^{x,z}_{f_X(x),f_Z(z)}\big)_{\mathsf{B}\mathsf{A}''}
  \ket{\hat\psi}\,,
  \tag{\path{eq:qld-combined-lines-consistency}}
\end{equation}
in which the point $(x,z,\alpha,\beta)$ is uniform on $\ell$ and the sum
ranges over pairs of line polynomials.  The proof then bounds $\Lambda$
twice, by different routes and with different results, and the printed
record is internally inconsistent: the display introducing $\Lambda$
asserts the first-route bound, while the closing lines assert the
second-route bound as the operative one; the lemma statement claims a
third form matching neither.

\paragraph{The first route.}
Three claims are chained.  The combined point measurement
$\hat Q^{x,z}$ is replaced by the ordered product of the two point
measurements at cost $\delta_Q^{1/2}$; the $X$-point factor is removed at
cost $m\,\delta_{\mathrm{Line}}^{1/2}$; and the remaining $Z$-point
correlation is shown to be within
$\sqrt{m}\,(\delta_P^{1/4} + \delta_Q^{1/4} + \eps^{1/4})$ of one.  The
factors of $m$ enter through the mixture structure of the sub-line
distribution: each restricted line distribution carries mixture weight
$1/(2m)$ in the line-point distribution, so an average bounded by
$\delta$ over the full distribution is bounded only by $2m\delta$ over a
restricted component and by $4m^2\delta$ over a product of two
components, and the Cauchy--Schwarz steps halve the exponents.  Chaining
the claims gives
\begin{equation*}
  \Lambda \,=\, O\big(m\,(\eps^{1/4} + \delta_P^{1/4} + \delta_Q^{1/4}
  + \delta_{\mathrm{Line}}^{1/2})\big)
  \,=\, m \cdot \poly(\eps,\, md/q)\,.
\end{equation*}
One intermediate assertion needs a harmless correction: the final line
of the proof of the second claim asserts the bound
$O(m^2 \delta_{\mathrm{Line}})$ for a consistency term that is close to
one, rather than for its deficit.  Read as a bound on the deficit ---
which is what the surrounding Cauchy--Schwarz argument requires and what
the invoked line--point consistency provides after the mixture
inflation --- the claim and its use are correct, with the square root of
the deficit giving the stated cost $m\,\delta_{\mathrm{Line}}^{1/2}$.

\paragraph{The second route.}
After the chain, the source restarts: it asserts that, by Property~2 of
its sub-line lemma, $\Lambda$ is a mixture, over types
$\mathrm{v}_1, \mathrm{v}_2 \in \{\mathrm{ALine}, \mathrm{DLine}\}$ and
indices $i, j \in \{1, \dots, m\}$, of terms
\begin{equation*}
  \E_{(\ell_X, x) \sim D_{\mathrm{v}_1,i}}\
  \E_{(\ell_Z, z) \sim D_{\mathrm{v}_2,j}}
  \sum_{f_X,f_Z} \bra{\hat\psi}
  \big(T^{\ell_X,\ell_Z}_{f_X,f_Z}\big)_{\mathsf{A}\mathsf{A}'} \ot
  \big(\Id - \hat{Q}^{x,z}_{f_X(x),f_Z(z)}\big)_{\mathsf{B}\mathsf{A}''}
  \ket{\hat\psi}
\end{equation*}
over two \emph{independently} sampled restricted line-point pairs,
bounds each term by $O(m^2\delta_P)$ using the restricted-distribution
consequence of the two-basis line lemma, and concludes
$\Lambda = O(m^2\delta_P)$.  The invoked decomposition concerns the
\emph{joint} law of the pair $((\ell_X, x), (\ell_Z, z))$, where $x$ and
$z$ are the two blocks of the single point sampled on $\ell$.  Property~2
of the sub-line lemma does not assert this joint law: it asserts the
component-wise product law of the pair of \emph{lines}
$(\ell_X, \ell_Z)$, and the conditional uniformity of each point on its
line separately --- equivalently, the laws of the two pairs
$((\ell_X, x), (\ell_Z, z'))$ and $((\ell_X, x'), (\ell_Z, z))$ in which
the \emph{other} point is resampled independently.  These marginal
statements are exactly what the first route's claims consume, each
involving one point at a time; the joint decomposition used by the
second route is a strictly stronger property and is not established by
the lemma it is attributed to.

\section{Why neither route yields the stated form}

The stated error form is $\poly(m^2\eps, md/q)$, a function of the two
quantities $m^2\eps$ and $md/q$ that vanishes as both vanish.  The
first-route bound is not dominated by any such function: along the
scaling $\eps = m^{-4}$ with $q/(md) \to \infty$, both $m^2\eps = m^{-2}$
and $md/q$ tend to $0$, while the term $m\,\eps^{1/4} = 1$ does not.
Hence the established derivation supports no bound of the claimed form.
The second-route bound, even granting the missing decomposition, carries
the separate prefactor $m^2$: writing $\delta_P = \poly(\eps, md/q)$
with some fixed positive exponents, the contribution
$m^2 (md/q)^{c}$ is not dominated by any vanishing function of $m^2\eps$
and $md/q$, as the scaling $\eps = 0$, $md/q = m^{-2/c}$ shows, and the
$\eps$-contribution fails similarly for sublinear exponents.  Neither
observation asserts that the \emph{conclusion} of the lemma is false for
some error function of the claimed form; the defect is that the printed
proof does not deliver it.

For orientation, the algebraic identity
$m\,\eps^{1/2} = (m^2\eps)^{1/2}$ shows what a derivation of the claimed
$\eps$-dependence would need: a chain that loses only square roots in the
$\eps$-type quantities \emph{after} the restricted-distribution
inflation would produce terms of the shape $(m^2\eps)^{1/2}$.  The
printed chain loses fourth roots, through two nested Cauchy--Schwarz
steps applied to already square-rooted quantities, and its
$md/q$-dependence enters as $m\,\delta_P^{1/4}$, which the claimed form
cannot absorb.

\section{Correction and alternative repairs}

The first direction below is the one the blueprint adopts; the remaining
directions are alternatives, none of which is claimed as established.

\paragraph{Corrected error form (adopted in the proof).}
The operative error form is
$\delta_{\mathrm{combine}} = m \cdot \poly(\eps, md/q)$ --- concretely,
the first-route bound.  Following the proof-gap protocol, the source-labelled
statement keeps its printed form and the correction is carried by the proof,
by the chapter remark, and by a separate formalization-support node, rather
than being substituted silently into the source statement.  The only consumer
is the polynomial-pair lemma, whose imported error
$a(md)^a(\delta_{\mathrm{combine}}^b + q^{-b} + 2^{-bmd})$ absorbs
$m^{O(1)}$ prefactors by enlarging $a$ and shrinking $b$.  The blueprint proof
records this absorption and the corrected deficit reading of the second claim.
Two obligations remain: carry out that absorption with explicit constants when
the lemma is formalized, and prove the three chained estimates.

\paragraph{Repairing the second route.}
Strengthen Property~2 of the sub-line lemma to the joint product law: in
each mixture component, the pairs $(\ell_X, x)$ and $(\ell_Z, z)$ are
independent with the stated restricted-distribution laws.  The sampling
procedure makes this plausible, since in each component the two pairs
are built from disjoint collections of the independent ingredients (the
two blocks of the uniform point, the direction blocks, and the fresh
auxiliary choices).  Open obligations: (i)~state and prove the
strengthened property for every component, including lines whose
direction blocks vanish and degenerate to singletons; (ii)~re-derive the
mixture decomposition of $\Lambda$ from the strengthened property;
(iii)~note that even then the route delivers only
$O(m^2\delta_P) = m^2 \cdot \poly(\eps, md/q)$, so this direction removes
the unsupported step of the printed proof but still requires the
error-form decision of the previous direction, with $m^2$ in place of
$m$.

\paragraph{Establishing the claimed form.}
Re-derive the chain so that every factor of $m$ enters jointly with
$\eps$ as $m^2\eps$, in the sense of the orientation identity above.
Open obligations: (i)~find a chaining that avoids one of the two nested
square-root losses, or prove restricted-distribution versions of the
underlying consistency lemmas directly at their inflated errors;
(ii)~handle the $md/q$-dependence, since a residual term
$m (md/q)^{c}$ is incompatible with the claimed form for the current
exponents.  No derivation along these lines is known here, and the
direction may be infeasible; in that case the corrected error form of
the first direction stands.

\section{Scalar and distributional scope of the formal claims}

The completed formal estimates corresponding to Claims~17-1 and~17-3
have a more limited scope than the source statements. Let
$D_{\mathrm{dir}}$ be a probability law on extended-line and source-line
triples whose extended-line marginal is the directly indexed law, with the
incidence and separate point-marginal properties of the sub-line lemma.
Write $\mathbb E_{\mathrm{dir}}$ for sampling this law and then a uniform
affine parameter on the extended line. For the point and line families
introduced above, define
\begin{align*}
 A_{\mathrm{dir}}&=\mathbb E_{\mathrm{dir}}\sum_{f_X,f_Z}
   \langle\hat\psi,(T_{f_X,f_Z}\otimes
     \hat Q^{x,z}_{f_X(x),f_Z(z)})\hat\psi\rangle,\\
 B_{\mathrm{dir}}&=\mathbb E_{\mathrm{dir}}\sum_{f_X,f_Z}
   \langle\hat\psi,(T_{f_X,f_Z}\otimes
     M^Z_{f_Z(z)}M^X_{f_X(x)})\hat\psi\rangle,\\
 Z_{\mathrm{dir}}&=\mathbb E_{\mathrm{dir}}\sum_{f_X,f_Z}
   \langle\hat\psi,(T_{f_X,f_Z}\otimes M^Z_{f_Z(z)})\hat\psi\rangle.
\end{align*}
Here $M^W$ abbreviates the expanded point effects at $x$ or $z$, with
placements $\mathsf{A}\mathsf{A}'$ and $\mathsf{B}\mathsf{A}''$ for the
two tensor factors. An undefined line evaluation selects the zero point
effect on the completed alphabet $\F_q\cup\{\bot\}$. The formal theorems give
\[
 |\Re A_{\mathrm{dir}}-\Re B_{\mathrm{dir}}|\le 2\sqrt{\delta_Q},
 \qquad
 |\Re Z_{\mathrm{dir}}-1|
 \le 2\sqrt m\,(\delta_P^{1/4}+\delta_Q^{1/4}+\eps^{1/4}).
\]

Claim~17-1 instead bounds the complex modulus $|A-B|$ for the source
law. The ordered product $M^Z_bM^X_a$ need not be Hermitian; a real-part
bound alone therefore does not prove the source assertion.
For $W_{a,b}=\hat Q_{a,b}-M^Z_bM^X_a$, its squared state norm is
$\langle\hat\psi,(I\otimes W_{a,b}^*W_{a,b})\hat\psi\rangle$.
The source prints $W_{a,b}^2$. The formal proof uses the adjoint square
and inserts the opposite-placed joint measurement in the distance
comparison, giving the same-placement bound $4\delta_Q$ and hence
$2\sqrt{\delta_Q}$. This constant is compatible with the source's
asymptotic error, but the real-part conclusion remains weaker.
For Claim~17-3 the opposite-placed positive effects commute, so
$Z_{\mathrm{dir}}$ is real; the distributional discrepancy remains.

The precise verdict is that these are proved auxiliary estimates, not
completed formalizations of the source claims. A source proof still needs
a complex-modulus comparison for Claim~17-1 and a transport from the directly
indexed law and completed evaluations to the source formulation.
The latter obligation includes the dimension-divisibility issue described
in \path{docs/paper-gaps/qpbt_ld-dimension-divisibility.tex}.
The source blueprint entries retain their statements without completion
marks; the real-part statements have separate formalization-support
entries. No hypothesis has been added to either Lean theorem.%
\footnote{Source: \path{claim:17-1}, lines~1140--1165, and
\path{claim:17-3}, lines~1204--1209, of
\path{references/qpbt-paper/14_analysis_of_the_pauli_basis_test.tex}.
The declarations are \leanid{subline\_replace\_by\_ordered\_product} and
\leanid{subline\_Z\_term\_near\_one} in
\path{MIPStarRE/QPBT/Combining/Claims.lean}; their separate blueprint
entries are \path{lem:claim-17-1-direct-real} and
\path{lem:claim-17-3-direct-real}.}

\section{Comparison with the blueprint}

The blueprint \cite{MIPStarREBlueprint} keeps the combined-lines lemma in
its source shape, error form $\poly(m^2\eps, md/q)$ included, and marks in
the proof that this form remains unproved: the proof derives the first route
only, records the established bound $m \cdot \poly(\eps, md/q)$, and uses
that bound in the sequel.  The
three claims are stated as separate lemmas with the corrected deficit
assertion for the second; the sub-line lemma states Property~2 at the
level of the two marginal laws and expressly disclaims the joint
conditional law of the two points, so the blueprint statement does not
license the second route; and the polynomial-pair lemma's proof absorbs
the $m$ prefactor into the constants of the imported soundness error.  A
chapter remark summarizes the discrepancy and cites this note.%
\footnote{The blueprint statements are \path{lem:qld-4-13},
\path{lem:qld-sublines}, \path{lem:claim-17-1}--\path{lem:claim-17-3},
\path{lem:qld-4-7}, and the remark
\path{rem:qld-4-13-source-defects}, all in
\path{blueprint/src/chapter/ch15_qpbt_combining.tex}.}

The source-labelled lemma links the carrier
\path{ExtendedLinesWitness} and the declaration
\path{exists_extendedLinesWitness}, which retains the printed error form
$\poly(m^2\eps,md/q)$ and whose proof is the tracked obligation.  The
established error form is linked from the separate formalization-support node
\path{lem:qld-4-13-established}, whose declaration
\path{exists_extendedLinesWitness_established} has the explicit form
$C m\,\poly(\eps,md/q)$ for a universal $C>0$; this agrees with the blueprint
convention that constants are absorbed into $\poly$.  Both proofs are open.

Neither declaration is source-shaped beyond its error form, and this note does
not claim otherwise.  Both return the carrier \path{ExtendedLinesWitness},
which replaces the source's lines in $\F_q^{2m+2}$ by the directly indexed
extended line space of
\path{docs/paper-gaps/qpbt_ld-dimension-divisibility.tex}, averages over the
corresponding directly indexed line-point law instead of over the line-point
distribution instantiated at dimension $2m+2$, and compares both sides of its
two consistency displays on the completed answer alphabet
$\F_q\cup\{\bot\}$, whose extra outcome collects the line polynomials that
admit no evaluation at the sampled point, in place of the source's answer
summation over $a\in\F_q$.  Neither blueprint node therefore carries a
statement-level formalization mark, and relating the two formulations is the
transport obligation of that note.

\section{Conclusion}

There is a proof-level gap in the cited lemma, together with an
intermediate inconsistency: the printed proof asserts two different
bounds for the same consistency defect, the first-route bound
$O(m(\eps^{1/4} + \delta_P^{1/4} + \delta_Q^{1/4} +
\delta_{\mathrm{Line}}^{1/2}))$ and the second-route bound
$O(m^2\delta_P)$, the second resting on a joint-law decomposition that
the sub-line lemma as stated does not supply; and the error form claimed
in the lemma statement, $\poly(m^2\eps, md/q)$, follows from neither.
The first route is sound after a harmless deficit correction in its
second claim and establishes $m \cdot \poly(\eps, md/q)$.  The
discrepancy is harmless downstream, since the sole consumer absorbs
polynomial-in-$m$ prefactors into its constants.  The blueprint therefore
retains the printed statement with the missing step explicitly tracked, uses
the corrected first-route form in the proof and in the sequel, and cites this
note; strengthening the sub-line lemma to license the second route is
optional, and no route to the printed form is currently known.  We do not
assert that the printed form is false.

\bibliographystyle{alpha}
\bibliography{references}

\end{document}
